\section{Preliminaries}
\label{sec:prelim}

We introduce notation and recall the building blocks used in our construction.
Throughout, $\mathbb{F}$ denotes a finite field of size $|\mathbb{F}| = p$ for a prime $p$, and $\lambda$ denotes the security parameter.
We write $\textsf{negl}(\lambda)$ for any function that is negligible in~$\lambda$, and $\ppt$ for probabilistic polynomial-time.

\subsection{Notation}
\label{sec:notation}

Matrices are denoted by bold uppercase letters $\mat{A}, \mat{B}, \mat{C} \in \mathbb{F}^{k \times k}$, and vectors by bold lowercase letters $\vect{x}, \vect{y}, \vect{z} \in \mathbb{F}^k$.
For a vector $\vect{x}$, we write $\vect{x}[i]$ for its $i$-th entry. For a
matrix $\mat{M}$, we write $\mat{M}[i,j]$ for the entry in the $i$-th row and $j$-th column. 
We denote the $i$-th row of $\mat{M}$ by $\mat{M}[i,*]$, and we will often write it simply as $\mat{M}[i]$. 
Similarly, we denote the $j$-th column of $\mat{M}$ by $\mat{M}[*,j]$. We write $[n] = \{0,1,\ldots,n-1\}$. For $t \in \mathbb{N}$, $I \sample [n]^t$ denotes the uniform sampling of $t$ indices from $[n]$. Similarly, $r \sample \mathbb{F}$ denotes uniform random sampling from
$\mathbb{F}$.



\subsection{Linear Error-Correcting Codes}
\label{sec:prelim-codes}

A \emph{linear code} over $\mathbb{F}$ is an $\mathbb{F}$-linear map
$\mathcal{C}: \mathbb{F}^k \to \mathbb{F}^n$ defined by a generator matrix
$\mat{G} \in \mathbb{F}^{k \times n}$, so that
$\mathcal{C}(\vect{m}) = \vect{m} \cdot \mat{G}$ for any input vector
$\vect{m} \in \mathbb{F}^k$.
The code has \emph{rate} $\rho = k/n$ and \emph{minimum distance} $d = \min_{\vect{c} \neq \vect{c}'} \Delta(\vect{c}, \vect{c}')$,
where the minimum is taken over distinct codewords
$\vect{c}, \vect{c}' \in \mathcal{C}$, and $\Delta$ denotes the Hamming distance.
The \emph{relative distance} is $\delta = d/n$.

We write $\enc(\vect{m}) \coloneqq \mathcal{C}(\vect{m})$ for the encoding of an input vector.
For a matrix $\mat{M} \in \mathbb{F}^{k \times k}$, we define the
\emph{row-wise encoding} $\widehat{\mat{M}} := \enc(\mat{M}) \in \mathbb{F}^{k \times n}$ as the matrix whose $i$-th row is $\enc(\mat{M}[i,*])$ for each $i \in [k]$.
The $j$-th column of $\widehat{\mat{M}}$ is denoted by
$\widehat{\mat{M}}[*,j] \in \mathbb{F}^k$ for each $j \in [n]$.

\begin{definition}[Interleaved code]\label{def:interleaved-code}
For an integer \(\ell\ge 1\), the \(\ell\)-fold interleaved code
\(\mathcal{C}^{\langle \ell\rangle}\) is obtained by grouping \(\ell\)
codewords of \(\mathcal{C}\) position by position.  More concretely, for
codewords
\[
  \vect{c}^{(0)},\ldots,\vect{c}^{(\ell-1)}\in\mathcal{C},
\]
their interleaving is the word
\(\vect{w}\in(\mathbb{F}^{\ell})^n\) whose \(j\)-th symbol is
\[
  \vect{w}[j]
  =
  \big(
    \vect{c}^{(0)}[j],
    \ldots,
    \vect{c}^{(\ell-1)}[j]
  \big)
  \in\mathbb{F}^{\ell}
  \quad\text{for }j\in[n].
\]
The code \(\mathcal{C}^{\langle \ell\rangle}\) is the set of all such
interleavings.

The row-wise encoding of a matrix is an interleaved codeword in this sense.
Indeed, for \(\mat{M}\in\mathbb{F}^{k\times k}\), the rows
\(\enc(\mat{M}[0,*]),\ldots,\enc(\mat{M}[k-1,*])\) are \(k\) codewords of
\(\mathcal{C}\), and their interleaving is exactly
\(\widehat{\mat{M}}\in\mathbb{F}^{k\times n}\).  The \(j\)-th interleaved
symbol is the encoded column
\(\widehat{\mat{M}}[*,j]\in\mathbb{F}^{k}\).
\end{definition}

The key property we exploit is \emph{linearity}: for any coefficients
$r_0, \ldots, r_{k-1} \in \mathbb{F}$,
\begin{equation}\label{eq:code-linearity}
\sum_{i=0}^{k-1} r_i \cdot \widehat{\mat{M}}[i,*]
=
\enc\!\left(\sum_{i=0}^{k-1} r_i \cdot \mat{M}[i,*]\right)
\end{equation}


\begin{definition}[Fold]\label{def:fold}
For vectors $\vect{c},\vect{r}\in \mathbb{F}^k$ with
$\vect{c}=(c_0,\ldots,c_{k-1})^\top$ and
$\vect{r}=(r_0,\ldots,r_{k-1})$, define
\[
\fold(\vect{r}, \vect{c}) \coloneqq \langle \vect{r}, \vect{c} \rangle = \sum_{i=0}^{k-1} r_i \cdot c_i.
\]
Equivalently, for the $j$-th column of a row-wise encoded matrix,
$\fold(\vect{r}, \widehat{\mat{M}}[*,j])$ equals the $j$-th coordinate of
$\enc(\vect{r}\mat{M})$.
This follows directly from the linearity of the code (Equation~\ref{eq:code-linearity}).
\end{definition}

\paragraph{Concrete instantiations}
Our protocol is compatible with any linear code.
Two natural choices are: (i) \emph{Reed--Solomon codes}, where $\mathcal{C}(\vect{m})$ evaluates the polynomial $\sum_i m_i X^i$ over a domain of size $n > k$, achieving the optimal (MDS) distance $\delta = 1 - k/n + 1/n$; and (ii) \emph{expander-based codes} as used in Brakedown~\cite{golovnev2023brakedown}, which achieve linear-time encoding $\mathcal{O}(n)$ with constant relative distance.


\subsection{Commitment Schemes}
\label{sec:prelim-commitments}
A commitment scheme allows a committer to bind itself to a message and later
open it. Some commitment schemes also hide the message before opening; our
soundness arguments use only binding. In this work, we use two types of
commitments: Merkle tree commitments and Pedersen vector commitments.

\subsubsection{Merkle tree commitments}
% A Merkle tree commitment over a collision-resistant hash function
% $H: \{0,1\}^* \to \{0,1\}^\lambda$ is a vector commitment for a vector
% $\vect{v}=(\vect{v}[1],\ldots,\vect{v}[n]) \in \mathbb{F}^n$.
% It consists of the following algorithms.

A Merkle tree commitment over a collision-resistant hash function $H: \{0,1\}^* \to \{0,1\}^\lambda$ is a vector commitment for a sequence $\vect{v}=(\vect{v}[j])_{j\in[n]}$.
The leaves may be randomized or deterministic. We write
\(\merkleC(\vect{v};\vect{o})\) for randomized leaves and
\(\merkleC(\vect{v};\bot)\) for deterministic leaves.

\begin{itemize}
\item $\rt \gets \merkleC(\vect{v};\vect{o})$:
outputs a root hash $\rt$ obtained by constructing a binary hash tree with
randomized leaves
% $
%   \ell_j = H(\vect{v}[j], \vect{o}[j])
%   \qquad \text{for } j\in[n].
% $
\[
\ell_j =
\begin{cases}
  H(\vect{v}[j], \vect{o}[j]) & \text{if } \vect{o}\neq\bot,\\
  H(\vect{v}[j]) & \text{if } \vect{o}=\bot.
\end{cases}
\]

Internal nodes are computed by hashing their two children.

\item $(\vect{v}[j],o_j,\pi_j) \gets \merkleO(\vect{v},\vect{o},j)$:
outputs the leaf value $\vect{v}[j]$, its opening randomness $o_j$, and its
authentication path $\pi_j$ for position $j\in[n]$.
In the deterministic case \(\vect{o}=\bot\), we omit the randomness and write
\((\vect{v}[j],\pi_j)\gets\merkleO(\vect{v},j)\).

\item $b \gets \merkleV(\rt,j,\vect{v}[j],\vect{o}[j],\pi_j)$:
recomputes the randomized leaf
$
  \ell_j = H(\vect{v}[j], \vect{o}[j])
$
and outputs $b=1$ if the authentication path is consistent with root $\rt$ at
position $j$, and outputs $b=0$ otherwise.
In the deterministic case, we write
\(\merkleV(\rt,j,\vect{v}[j],\pi_j)\).
\end{itemize}


Merkle tree commitments are \emph{position-binding}: under the collision
resistance of $H$, no $\ppt$ adversary can produce a root $\rt$, an index
$j\in[n]$, two distinct openings $(v,o)\neq(v',o')$, and authentication paths
$\pi_j,\pi'_j$ such that
\[
  \merkleV(\rt,j,v,o,\pi_j)
  =
  \merkleV(\rt,j,v',o',\pi'_j)
  =
  1,
\]
except with probability
$\epsilon_{\mathsf{merkle}}(\lambda)=\mathsf{negl}(\lambda)$.



\subsubsection{Pedersen commitments}
Pedersen commitments can be viewed as algebraic vector commitments over a group $\mathbb{G}$ of prime order $q$. Let $\mathsf{ck}=(G_1,\ldots,G_n,H)$ be a commitment key consisting of fixed generators in $\mathbb{G}$. A Pedersen commitment is defined by the following algorithms.

\begin{itemize}
\item $\cm \gets \pedC(\ck,\vect{v};o)$:
outputs a commitment $\cm := \sum_{i=1}^{n}\vect{v}[i]G_i + oH \in \mathbb{G}$ for a vector $\vect{v}\in\mathbb{F}^n$ and randomness $o\sample\mathbb{F}$.

\item $\pi \gets \pedO(\vect{v},o)$:
outputs the opening $\pi=(\vect{v},o)$ consisting of the committed vector and the randomness.

\item $b \gets \pedV(\ck,\cm,\vect{v},o)$:
outputs $b=1$ if $\cm \stackrel{?}{=} \sum_{i=1}^{n}\vect{v}[i]G_i + oH$,
and outputs $b=0$ otherwise.
\end{itemize}

Pedersen commitments are \emph{perfectly hiding} and
\emph{computationally binding}: under the discrete logarithm assumption,
no $\ppt$ adversary can produce a commitment $\cm$, two distinct vectors
$\vect{v}\neq\vect{v}'$, and randomness values $o,o'$ such that
\[
\pedV(\ck,\cm,\vect{v},o) = \pedV(\ck,\cm,\vect{v}',o') = 1,
\]
except with probability
$\epsilon_{\mathsf{ped}}(\lambda)=\mathsf{negl}(\lambda)$.

They are also \emph{additively homomorphic}: for
$\cm_1=\pedC(\ck,\vect{v}_1;o_1)$ and
$\cm_2=\pedC(\ck,\vect{v}_2;o_2)$, $\cm_1+\cm_2 = \pedC(\ck,\vect{v}_1+\vect{v}_2;o_1+o_2).$


\subsection{Succinct Non-interactive Arguments of Knowledge}
\label{sec:prelim-snark}

A SNARK (succinct non-interactive argument of knowledge) for an NP relation
$\relation$ is a tuple of algorithms $\bPi=(\setup,\prove,\verify)$.

\begin{itemize}
\item $\para \gets \bPi.\setup(1^\lambda,\relation)$:
generates public parameters $\para$ for the relation $\relation$.

\item $\pi \gets \bPi.\prove(\para,\stmt,\wit)$:
outputs a proof $\pi$ for a statement $\stmt$ and witness $\wit$ satisfying $(\stmt,\wit)\in\relation$.

\item $b \gets \bPi.\verify(\para,\stmt,\pi)$:
outputs $b=1$ if $\pi$ is accepted as a valid proof for $\stmt$, and outputs $b=0$ otherwise.
\end{itemize}

\begin{remark}[Security Properties]
A SNARK scheme $\bPi$ satisfies completeness, if an honestly generated proof for any valid statement-witness pair $(\stmt,\wit)\in\relation$ is accepted by the verifier with overwhelming probability. It satisfies knowledge soundness if any PPT prover that produces an accepting proof for a statement $\stmt$ must know a valid witness $\wit$ such that $(\stmt,\wit)\in\relation$, except with negligible probability. This is formalized by the existence of an efficient knowledge extractor. A SNARK may additionally satisfy zero knowledge, meaning that the proof reveals no information beyond the truth of the statement. This is modeled by a simulator $\mathsf{Sim}$ that can generate a proof indistinguishable from a real proof without knowing a valid witness. A SNARK satisfying zero knowledge is called a zk-SNARK.
\end{remark}


\subsubsection{Commit-and-Prove SNARKs}
A commit-and-prove SNARK (CP-SNARK)~\cite{gnark, legosnark} is a variant of a SNARK that proves statements about witnesses bound to an external commitment. We call such a commitment a \emph{witness commitment}. Let $\com$ be a commitment scheme with commitment key $\ck$. For an NP relation $\relation$ over tuples $(\stmt,\comwit,\omega)$, a CP-SNARK proves knowledge of a witness
$\wit := (\comwit,\omega)$ such that $(\stmt,\comwit,\omega)\in\relation$,
where $\comwit$ is the committed part of the witness and $\omega$ is the
remaining private witness.  The committed witness is bound separately by a
witness commitment
$\widetilde{\cm}=\com(\ck,\comwit;\tilde{o})$, where $\tilde{o}$ is the
opening randomness.  Equivalently, the CP proof is verified with respect to the
extended public statement $\stmt_{\cp}:=(\stmt,\widetilde{\cm})$.  When the
context is clear, we write this extended statement simply as $\stmt$.
A CP-SNARK scheme is a tuple of algorithms
$\bPi_{\cp}=(\setup,\com,\prove,\verify)$.

\begin{itemize}
    \item $\para \gets \bPi_{\cp}.\setup(1^\lambda,\relation)$: generates public parameters $\para=(\pk,\vk,\ck)$ for the relation $\relation$, where $\pk$ is the proving key, $\vk$ is the verifying key, and $\ck$ is the commitment key.

    \item $\widetilde{\cm} \gets \bPi_{\cp}.\com(\para,\comwit;\tilde{o})$:
    outputs a witness commitment to the committed witness part $\comwit$.

    \item $\pi_{\cp} \gets \bPi_{\cp}.\prove(\para,\stmt_{\cp},\wit)$:
    outputs a proof for the extended statement $\stmt_{\cp}$ and witness
    $\wit=(\comwit,\omega)$.

    \item $b \gets \bPi_{\cp}.\verify(\para,\stmt_{\cp},\pi_{\cp})$:
    outputs $b=1$ if $\pi_{\cp}$ is accepted as a valid proof for the extended
    statement $\stmt_{\cp}$, and outputs $b=0$ otherwise. Equivalently, we may
    write verification as
    $\bPi_{\cp}.\verify(\para,\stmt,\widetilde{\cm},\pi_{\cp})$.
\end{itemize}




\subsection{CP-Link}
\label{sec:prelim-cplink}

A CP-Link proof is a proof system for linking commitments. In its basic
equality-link form, it proves that two or more commitments are commitments to
the same message, without revealing the message or the opening randomness.

Let $\com$ be a commitment scheme. For commitment keys
$\ck_1,\ldots,\ck_t$ and commitments $\cm_1,\ldots,\cm_t$, the linking relation is defined as
\[
\begin{aligned}
\relation_{\mathsf{link}}
=
\Big\{
&\big(
((\ck_i,\cm_i)_{i=1}^{t}),
(m,(o_i)_{i=1}^{t})
\big)
:\\
&\cm_i = \com(\ck_i,m;o_i)
\quad \text{for all } i\in[t]
\Big\}.
\end{aligned}
\]
A proof for $\relation_{\mathsf{link}}$ convinces the verifier that all commitments $\cm_1,\ldots,\cm_t$ are commitments to the same message $m$.

We use CP-Link through the standard statement/witness interface for a linking
relation.  For the equality-link relation above, the public statement and
witness are
\[
  \stmt_{\mathsf{link}}
  =
  ((\ck_i,\cm_i)_{i=1}^{t}),
  \qquad
  \wit_{\mathsf{link}}
  =
  (m,(o_i)_{i=1}^{t}).
\]
More structured linking relations may link commitments to deterministic blocks
or projections of a common message; in that case the projection maps are part of
\(\relation_{\mathsf{link}}\), but the algorithmic interface is unchanged.

A CP-Link proof system is a tuple of algorithms
$\bPi_{\mathsf{link}}=(\setup,\prove,\verify)$.

\begin{itemize}
\item $\para_{\mathsf{link}} \gets
\bPi_{\mathsf{link}}.\setup(1^\lambda,\relation_{\mathsf{link}})$:
generates public parameters for the linking relation.

\item $\pi_{\mathsf{link}} \gets
\bPi_{\mathsf{link}}.\prove(
\para_{\mathsf{link}},
\stmt_{\mathsf{link}},
\wit_{\mathsf{link}})$:
outputs a proof that the statement-witness pair satisfies the linking relation.

\item $b \gets
\bPi_{\mathsf{link}}.\verify(
\para_{\mathsf{link}},
\stmt_{\mathsf{link}},
\pi_{\mathsf{link}})$:
outputs $b=1$ if the proof is accepted, and outputs $b=0$ otherwise.
\end{itemize}

\subsection{Freivalds' Algorithm}
\label{sec:prelim-freivalds}

Freivalds' algorithm~\cite{freivalds1977probabilistic} is a classical randomized algorithm for verifying matrix multiplication.
Given $\mat{A}, \mat{B}, \mat{C} \in \mathbb{F}^{k \times k}$, it samples $\vect{r} \sample \mathbb{F}^k$ and checks whether $\vect{r}(\mat{A}\mat{B}) = \vect{r}\mat{C}$.
The correctness relies on the following:

\begin{lemma}[Schwartz--Zippel~\cite{schwartz1980fast,zippel1979probabilistic}]\label{lem:sz}
Let $P(X_1, \ldots, X_m)$ be a non-zero polynomial of total degree $d$ over a finite field $\mathbb{F}$.
Then for $r_1, \ldots, r_m \sample \mathbb{F}$ chosen uniformly and independently,
\[
\Pr[P(r_1, \ldots, r_m) = 0] \le \frac{d}{|\mathbb{F}|}.
\]
\end{lemma}

In our protocol, we use a structured variant: instead of sampling $\vect{r} \sample \mathbb{F}^k$ uniformly, we sample one scalar $r \sample \mathbb{F}$ and set $\vect{r} = (1, r, r^2, \ldots, r^{k-1})$.
This keeps the Freivalds challenge compact in the transcript and in the SNARK
statement: the verifier records one field element rather than $k$ independent
ones. If $\mat{D} = \mat{A}\mat{B} - \mat{C} \neq \mat{0}$, then for any non-zero
column $\vect{d}_j$ of $\mat{D}$, the value $\langle \vect{r}, \vect{d}_j \rangle$ is a non-zero univariate polynomial in
$r$ of degree at most $k-1$.
By Lemma~\ref{lem:sz}, it vanishes with probability at most $(k-1)/|\mathbb{F}|$.
